\TieudeT{PHONG TỎA VẬT LÍ 12 - VẬT LÍ NHIỆT}
\subsubsection*{PHẦN I. Thí sinh trả lời câu hỏi từ câu 1 đến câu 18. Mỗi câu hỏi thí sinh chỉ chọn một phương án}
\setcounter{ex}{0}
\begin{ex} %Câu 1
	Nhiệt độ của nước đá đang tan theo thang nhiệt độ Celsius là
	\choice
	{$100^\circ\text{C}$}
	{\True $0^\circ\text{C}$}
	{$273K$}
	{$373K$}
	\loigiai{}
\end{ex}

\begin{ex} %Câu 2
	Trong thời gian sắt đông đặc, nhiệt độ của nó
	\choice
	{\True không đổi}
	{không ngừng giảm}
	{mới đầu tăng, sau giảm}
	{không ngừng tăng}
	\loigiai{}
\end{ex}

\begin{ex} %Câu 3
	Trong quá trình chất khí nhận nhiệt lượng và nhận công thì $Q$ và $A$ trong biểu thức $\Delta U = Q + A$ phải thoả mãn điều kiện
	\choice
	{$Q < 0$, $A > 0$}
	{$Q > 0$, $A < 0$}
	{\True $Q > 0$, $A > 0$}
	{$Q < 0$, $A < 0$}
	\loigiai{}
\end{ex}

\begin{ex} %Câu 4
	Nhiệt năng là
	\choice
	{động năng chuyển động của phân tử}
	{động năng chuyển động của vật}
	{\True tổng động năng của các phân tử cấu tạo nên vật}
	{tổng thế năng và động năng}
	\loigiai{}
\end{ex}

\begin{ex} %Câu 5
Phát biểu nào sau đây là đúng?
	\choice
	{Độ biến thiên nội năng của một vật là độ biến thiên nhiệt độ của vật đó}
	{Nội năng gọi là nhiệt lượng}
	{Nội năng là phần năng lượng vật nhận được hay mất bớt đi trong quá trình chuyển động}
	{\True Có thể làm thay đổi nội năng của vật bằng cách thực hiện công và truyền nhiệt}
	\loigiai{}
\end{ex}

\begin{ex} %Câu 6
	Chất rắn được phân loại thành
	\choice
	{chất rắn đơn tinh thể và chất rắn vô định hình}
	{\True chất rắn kết tinh và chất rắn vô định hình}
	{chất rắn đa tinh thể và chất rắn vô định hình}
	{chất rắn đơn tinh thể và chất rắn đa tinh thể}
	\loigiai{}
\end{ex}

\begin{ex} %Câu 7
Phát biểu nào sau đây là \textbf{sai}?
	\choice
	{Khoảng cách giữa các phân tử trong chất lỏng nhỏ hơn trong chất khí}
	{Khoảng cách giữa các phân tử trong chất rắn nhỏ hơn trong chất khí}
	{\True Khoảng cách giữa các phân tử trong chất rắn lớn hơn trong chất khí}
	{Khoảng cách giữa các phân tử trong chất khí có thể thay đổi}
	\loigiai{}
\end{ex}

\begin{ex} %Câu 8
	Trong sự dẫn nhiệt, nhiệt tự truyền
	\choice
	{từ vật có nhiệt năng lớn hơn sang vật có nhiệt năng nhỏ hơn}
	{từ vật có khối lượng lớn hơn sang vật có khối lượng nhỏ hơn}
	{\True từ vật có nhiệt độ cao hơn sang vật có nhiệt độ thấp hơn}
	{từ vật có nhiệt độ thấp sang nhiệt độ cao hơn}
	\loigiai{}
\end{ex}

\begin{ex} %Câu 9
	Đặc điểm nào sau đây là đặc điểm của thể lỏng?
	\choice
	{Khoảng cách giữa các phân tử rất lớn so với kích thước của chúng}
	{\True Lực tương tác phân tử yếu hơn lực tương tác phân tử ở thể rắn}
	{Không có thể tích và hình dạng riêng xác định}
	{Các phân tử dao động xung quanh vị trí cân bằng xác định}
	\loigiai{}
\end{ex}

\begin{ex} %Câu 10
Phát biểu nào sau đây là \textbf{sai}. Nhiệt nóng chảy riêng của chất rắn
	\choice
	{là nhiệt lượng cần cung cấp để làm nóng chảy hoàn toàn $1\ kg$ chất đó ở nhiệt độ nóng chảy}
	{\True là nhiệt lượng cần cung cấp để làm nóng chảy hoàn toàn lượng chất đó ở nhiệt nóng chảy}
	{phụ thuộc vào bản chất của chất rắn}
	{có đơn vị là $J/kg$}
	\loigiai{}
\end{ex}


\begin{ex}%Câu 11
	Cách nào sau đây \textbf{không} làm thay đổi nhiệt năng của vật?
	\choice
	{Cọ xát vật với một vật khác}
	{Đốt nóng vật}
	{Cho vật vào môi trường có nhiệt độ thấp hơn vật}
	{\True Treo vật lên giá đỡ}
	\loigiai{
		
	}
\end{ex}

\begin{ex}%Câu 12
	Nhiệt kế nào sau đây có thể đo được nhiệt độ của nước đang sôi?
	\choice
	{\True Nhiệt kế thủy ngân}
	{Nhiệt kế rượu}
	{Nhiệt kế y tế}
	{Nhiệt kế rượu và nhiệt kế thủy ngân}
	\loigiai{
		
	}
\end{ex}

\begin{ex}%Câu 13
	
	Một thang đo X lấy điểm băng là $-10X$, lấy điểm sôi là $90X$. Nhiệt độ của một vật đọc được trên nhiệt kế Celsius là $40^\circ C$ thì trên nhiệt kế X có nhiệt độ bằng
	\choice
	{$20\ X$}
	{\True $30\ X$}
	{$40\ X$}
	{$50\ X$}
	\loigiai{
		\begin{itemize}
			\item $\dfrac{40-0}{100-0}=\dfrac{t - (-10X)}{90X-(-10X)} \Rightarrow t = 30\ X$.
		\end{itemize}
	}
\end{ex}

\begin{ex}%Câu 14
	Nhiệt độ ở khu vực phía trên bề mặt nước biển bị ảnh hưởng mạnh bởi nhiệt dung riêng của nước. Một trong những nguyên nhân là khi $1\ m^3$ nước nguội đi $1^\circ C$ nó sẽ làm cho nhiệt độ của một thể tích không khí lớn tăng lên $1^\circ C$. Biết nhiệt dung riêng của không khí và của nước xấp xỉ bằng $1000\ J/kg.K$ và $4186\ J/kg.K$; khối lượng riêng của không khí và của nước lần lượt là $1,3\ kg/m^3$ và $1000\ kg/m^3$. Thể tích của khối không khí này \textbf{xấp xỉ} bằng
	\choice
	{$1220\ m^3$}
	{$2230\ m^3$}
	{\True $3220\ m^3$}
	{$2203\ m^3$}
	\loigiai{
		\begin{itemize}
			\item  $Q_{thu} = Q_{\text{tỏa}} \Rightarrow m_k.c_k.\Delta t_k = m_n.c_n.\Delta t_n \Rightarrow D_kV_k.c_k = D_n.V_n.c_n \Rightarrow V_k = \dfrac{D_n.V_n.c_n}{D_k.c_k} = \dfrac{1000.1.4186}{1,3.1000} = 3220\ m^3$.
		\end{itemize}
	}
\end{ex}

\begin{ex}%Câu 15
	Một thước $cm$ được đặt dọc theo một nhiệt kế thủy ngân chưa được chia vạch như hình dưới đây. Trên nhiệt kế chỉ đánh dấu điểm đóng băng và điểm sôi của nước tinh khiết ở áp suất tiêu chuẩn. Giá trị nhiệt độ đang hiển thị trên nhiệt kế là
	\begin{center}
		\begin{tikzpicture}[draw=Mapcolor, very thick, >=stealth']%nhiệt kế thủy ngân
			%\draw[dotted] (-2,-1) grid (14,2);
			\draw [rounded corners=0.2cm] (0.4,-0.25) rectangle (13,0.25);
			\draw [rounded corners=0.1cm] (-0.4,-0.1) rectangle (12.2,0.1);
			\draw [rounded corners=0.1cm, fill=Mapcolor] (-0.4,-0.1) rectangle (6.6,0.1);
			\foreach \x in {0,0.1,...,12}{
				\draw (\x,-0.3) -- (\x, -0.5);
			}
			\foreach \x in {0,...,12}{
				\draw (\x,-0.3) -- (\x,-0.7) node [below] {\x};
			}
			\draw (1,-0.3) -- (1,1) node [above] {Điểm đóng băng};
			\draw (11,-0.3) -- (11,1) node [above] {Điểm sôi};
			\draw (6.6,0) -- (6.6,1) node[above] {$6,6$};
		\end{tikzpicture}
	\end{center}
	\choice
	{$56^{\circ}C$}
	{$60^{\circ}C$}
	{$66^{\circ}C$}
	{$44^{\circ}C$}
	\loigiai{
		\begin{itemize}
			\item Ta có: $\dfrac{\Delta x}{\Delta t\left( ^\circ C\right)}  = \text{hằng số} \Rightarrow \dfrac{11-1}{100 - 0} = \dfrac{6,6-1}{t-0} \Rightarrow t \approx 56^\circ C$.
		\end{itemize}
		
	}
\end{ex}

\begin{ex}%Câu 16
	Một ấm nhôm có khối lượng $m_b = 600\ g$ chứa $V = 1,5$ lít nước ở $t_1 = 20^\circ C$, sau đó đun bằng bếp điện. Sau $t = 35$ phút thì đã có $20\%$ khối lượng nước đã hóa hơi ở nhiệt độ sôi $t_2 = 100^\circ C$. Biết rằng, $H = 75\%$ nhiệt lượng mà bếp cung cấp được dùng vào việc đun nước. Cho biết nhiệt dung riêng của nước là $c_n = 4190\ J/kg.K$, của nhôm là $c_b = 880\ J/kg.K$, nhiệt hóa hơi riêng của nước ở $100^\circ C$ là $L = 2,26.10^6\ J/kg$, khối lượng riêng của nước là $D = 1\ kg/\text{lít}$. Công suất cung cấp nhiệt của bếp điện \textbf{gần giá trị nào nhất} sau đây?
	\choice
	{\True $776\ W$}
	{$796\ W$}
	{$786\ W$}
	{$876\ W$}
	\loigiai{
		\begin{itemize}
			\item $Q = \left( m_b.c_b+m_n.c_n\right) \Delta t + 0,2.m_n.L \Rightarrow \mathcal{H.P}.t = \left( m_b.c_b+m_n.c_n\right) \Delta t + 0,2.m_n.L \Rightarrow 0,75.\mathcal{P}.35.60 = \left(0,6.880 + 1,5.4190 \right).80 + 0,2.1,5.2,26.10^6 \Rightarrow \mathcal{P} \approx 776,5\ W$. 
		\end{itemize}
	}
\end{ex}

\begin{ex}%Câu 17
\impicinpar{Để kiểm tra thời gian ngắt mạch của một cầu chì khi xảy ra đoản mạch, một học sinh mắc cầu chì vào một nguồn điện như hình vẽ. Chì có nhiệt dung riêng là $130\ J/kg.K$, nhiệt độ nóng chảy là $327,5\ ^\circ C$. Nguồn điện có suất điện động $\mathcal{E} = 12\ V$ và điện trở trong $r = 0,25\ \Omega$. Cho rằng cầu chì sẽ đứt ngay khi đạt nhiệt độ nóng chảy. Cho rằng nhiệt độ ban đầu của dây chì bằng nhiệt độ phòng là $27,5\ ^\circ C$. Dây chì có điện trở $R = 11,75\ \Omega$, và khối lượng $m = 0,1\ g$. Thời điểm mạch bị ngắt bởi cầu chì kể từ thời điểm đóng mạch là}{\begin{tikzpicture}[draw=Mapcolor, very thick, >=latex']
			%	\draw [dotted] (-2,-2) grid (2,2);
			\draw (-0.15,-1) -- (-2,-1) -- (-2,1) -- (2,1) -- (2,0.25)
			(2,-0.25) -- (2,-1) -- (0.15,-1);
			\draw [fill = Mapcolor] (-0.75,0.8) rectangle (0.75,1.2);
			%	\fill [white] (-0.15,-1.1) rectangle (0.15,-0.9);
			\draw (-0.15,-0.8) -- (-0.15,-1.2);
			\draw (0.15,-0.9) -- (0.15,-1.1);
			\fill [white] (1.8,-0.25) -- (2.2,0.25);
			\draw (2,-0.25) -- (2.3,0.15) node [right] {K};
			\draw node [below] at (0,-1.2) {$\mathcal{E}, r$}
			node [above] at (0,1.25) {Cầu chì};
			
			
	\end{tikzpicture}}
	
	\choice
	{$0,25\ s$}
	{\True $0,33\ s$}
	{$0,38\ s$}
	{$0,16\ s$}
	\loigiai{
		\begin{itemize}
			\item Cường độ dòng điện qua mạch: $ I = \dfrac{\mathcal{E}}{R+r} = \dfrac{12}{11,75+0,25} = 1\ A$.
			\item Thời điểm ngắt mạch: $P = \dfrac{A}{t} \Rightarrow t = \dfrac{A}{P} = \dfrac{m.c\Delta t}{I^2.R} = \dfrac{0,1.10^{-3}.130.300}{1^2.11,75} \approx 0,33\ s$.
		\end{itemize}
	}
\end{ex}
\begin{ex}%Câu 18
	Có hai quả cầu bằng chì giống nhau cùng khối lượng có nhiệt dung riêng là $c$, chuyển động đến va chạm mềm trực diện với tốc độ lần lượt là $v$ và $2v$. Giả thiết toàn bộ phần năng lượng tiêu hao do va chạm được hai quả cầu hấp thụ hoàn toàn. Độ tăng nhiệt độ $\Delta t$ của hai quả cầu là
	\choice
	{\True $\dfrac{9v^2}{8c}$}
	{$\dfrac{7v^2}{8c}$}
	{$\dfrac{9v^2}{7c}$}
	{$\dfrac{11v^2}{7c}$}
	\loigiai{
		\begin{itemize}
			\item Bảo toàn động lượng: $\overrightarrow{p_1}+\overrightarrow{p_2} = \overrightarrow{p} \Rightarrow \left| p_1 - p_2\right| = p \Rightarrow \left| mv - 2mv\right|  = (m+m)V \Rightarrow V = \dfrac{1}{2}v$.
			\item $Q = 2mc\Delta t \Rightarrow W_{\text{đ1}} + W_{\text{đ2}} - W_{\text{đ}} = 2mc\Delta t \Rightarrow \dfrac{1}{2}v^2 + \dfrac{1}{2}(2v)^2 - \dfrac{1}{2}2.\left( \dfrac{1}{2}v\right)^2 = 2c\Delta t \Rightarrow \Delta t = \dfrac{9v^2}{8c} $.
			
		\end{itemize}
	}
\end{ex}
\subsubsection*{PHẦN 2. Thí sinh trả lời từ câu 1 đến câu 4. Trong mỗi ý a), b), c), d) ở mỗi câu, thí sinh chọn đúng hoặc sai.}
\setcounter{ex}{0}
\begin{ex}
	Một học sinh pha chế một mẫu trà sữa bằng cách trộn các mẫu chất lỏng với nhau: nước trà đen (mẫu A), nước đường nâu (mẫu B) và sữa tươi (mẫu C). Các mẫu chất lỏng này chỉ trao đổi nhiệt lẫn nhau mà không gây ra các phản ứng hoá học. Bỏ qua sự trao đổi nhiệt với môi trường và bình chứa. Nhiệt độ trước khi trộn của mẫu A, mẫu B và mẫu C lần lượt là $10^{\circ}C$, $15^{\circ}C$ và $20^{\circ}C$. Biết rằng:
	- Khi trộn mẫu A với mẫu B với nhau thì nhiệt độ cân bằng của hệ là $13^{\circ}C$.\\
	- Khi trộn mẫu B với mẫu C với nhau thì nhiệt độ cân bằng của hệ là $18^{\circ}C$.
	\choiceTFt
	{\True Khi trộn mẫu A với mẫu B với nhau thì sau khi đạt trạng thái cân bằng nhiệt thì nhiệt độ của mẫu B giảm đi $2\ K$}
	{Nhiệt độ cân bằng của hệ khi trộn mẫu A với mẫu C là $16^{\circ}C$}
	{Nhiệt độ cân bằng của hệ khi trộn cả ba mẫu với nhau là $15,5^{\circ}C$}
	{\True Nếu học sinh này pha thêm một mẫu sữa tươi nữa vào hỗn hợp ba mẫu ở câu c thì nhiệt độ cân bằng của hệ lúc này là $17,5^{\circ}C$}
	\loigiai{
		Gọi $C_A$, $C_B$, $C_C$ lần lượt là nhiệt dung của $A$, $B$, $C$.
		\begin{itemchoice}
			\itemch Nhiệt độ mẫu $B$ giảm là $15-13 = 2^\circ C$.
			\itemch
			\begin{itemize}
				\item $\begin{cases}
					A+B: C_A(13-10) &= C_B(15-13)\\
					B+C: C_B(18-15) &= C_C(20-18)
				\end{cases} \Rightarrow 
				\begin{cases}
					C_A &= \dfrac{2}{3}C_B\\
					C_C &= \dfrac{3}{2}C_B
				\end{cases} 
				$
				. Chọn $C_B=1 \Rightarrow C_A = \dfrac{2}{3}$ và $C_C = \dfrac{3}{2}$.
				\item Khi trộn $A$ và $C$: $C_A(t-10) = C_C (20-t) \Rightarrow \dfrac{2}{3}(t-10) = \dfrac{3}{2}(20-t) \Rightarrow t \approx 16,9^\circ C$.
			\end{itemize}
			
			\itemch Khi trộn 3 mẫu: $C_A(t_A - t_{cb}) + C_B(t_B - t_{cb}) + C_C(t_C - t_{cb}) = 0 \Rightarrow t_{cb} \approx 16,3^\circ C$.
			\itemch $C_A(t_A - t_{cb}) + C_B(t_B - t_{cb}) + 2.C_C(t_C - t_{cb}) = 0 \Rightarrow t_{cb} = 17,5^\circ C$.
		\end{itemchoice}
	}
\end{ex}

\begin{ex}
	Khi cung cấp nhiệt lượng $2\ J$ cho khí trong xilanh đặt nằm ngang, khí nở ra đẩy pit-tông di chuyển đều đi được $5\ cm$. Cho lực ma sát giữa pittông và xilanh là $10\ N$. Bỏ qua áp suất khí quyển bên ngoài pittông.
	\choiceTFt
	{\True Quá trình trên hệ nhận nhiệt lượng nên $Q > 0$}
	{Độ lớn của công chất khí thực hiện để pittông chuyển động đều là $5\ J$}
	{\True Quá trình trên khí thực hiện công nên $A < 0$}
	{Độ biến thiên nội năng của khí là $15\ J$}
	\loigiai{
		\begin{itemchoice}
			\itemch $Q = 2\ J$.
			\itemch Độ lớn công của chất khí thực hiện: $\left| A\right| = F_{ms}.s = 10.0,05 = 0,5\ J$.
			\itemch $A = -0,5\ J$.
			\itemch $\Delta U = Q + A = 2 + (-0,5) = 1,5\ J$.
		\end{itemchoice}
	}
\end{ex}

\begin{ex}
	Một nhà máy thép mỗi lần luyện được $35$ tấn thép. Cho nhiệt nóng chảy riêng của thép là $2,77.10^5\ J/kg$. Giả sử nhà máy sử dụng khí đốt để nấu chảy thép trong lò thổi (nồi nấu thép). Biết khi đốt cháy hoàn toàn $1\ kg$ khí đốt thì nhiệt lượng tỏa ra là $44.10^6\ J$.
	\choiceTFt
	{\True Nhiệt lượng cần cung cấp để làm nóng chảy thép trong mỗi lần luyện của nhà máy ở nhiệt độ nóng chảy là $9,7\ GJ$}
	{Nếu có $80\%$ nhiệt lượng khí tỏa ra dùng để làm nóng chảy thép ở nhiệt độ nóng chảy thì khối lượng khí đốt cần sử dụng là $220\ kg$}
	{\True Nếu có $80\%$ nhiệt lượng khí tỏa ra dùng để làm nóng chảy thép ở nhiệt độ nóng chảy và thời gian mỗi lần luyện là $30$ phút thì trong $1$ phút khối lượng khí cần đốt là $9,2\ kg$}
	{\True Việc sử dụng khí đốt để vận hành các nhà máy thép có thể gây ô nhiễm môi trường, gây nguy hiểm đối với con người và sinh vật}
	\loigiai{
		\begin{itemchoice}
			\itemch $Q_{nc} = m.\lambda = 35.10^3.2,77.10^5 \approx 9,7.10^9\ J = 9,7\ GJ$.
			\itemch $Q = \dfrac{Q_{nc}}{\mathcal{H}} \Rightarrow m.q = \dfrac{Q_{nc}}{\mathcal{H}} \Rightarrow m \approx 275\ kg$.
			\itemch Khối lượng khí đốt trong 1 phút là $\dfrac{m}{t} = \dfrac{275}{30} \approx 9,2\ kg$.
			\itemch 
		\end{itemchoice}
	}
\end{ex}

\begin{ex}
	Một ô tô gia đình có trọng lượng $1,5.10^4\ N$ chạy thẳng với tốc độ không đổi $72\ km/h$, trong một giờ tiêu thụ $8\ l$ xăng. Biết ô tô chịu lực cản có độ lớn bằng $0,08$ lần trọng lượng của nó, khối lượng riêng và năng suất tỏa nhiệt của xăng lần lượt là $700\ kg/m^3$ và $q = 4,6. 10^7\ J/kg$.
	\choiceTFt
	{\True Quãng đường ô tô đã đi được là $72\ km$}
	{Nhiệt lượng tỏa ra khi đốt cháy hoàn toàn xăng là $2,576. 10^7\ J$}
	{\True Công suất lực kéo của ô tô là $24\ kW$}
	{\True Hiệu suất của động cơ ô tô xấp xỉ $33,54\%$}
	\loigiai{
		\begin{itemchoice}
			\itemch $s = v.t = 72.1 = 72\ km$.
			\itemch $Q = q.m = q.D.V = 4,6.10^7.8.10^{-3}.700 = 25,76.10^7\ J$.
			\itemch $\mathcal{P} = F.v \Rightarrow \mathcal{P} = 0,08P.v = 0,08.1,5.10^4.20 = 24000\ W = 24\ kW$.
			\itemch $\mathcal{H} = \dfrac{A}{Q} = \dfrac{F.s}{Q} = \dfrac{0,08.1,5.10^4.72000}{25,76.10^7} \approx 0,3354 = 33,54\%$.
		\end{itemchoice}
	}
\end{ex}
\subsubsection*{PHẦN 3. Thí sinh trả lời từ câu 1 đến câu 6.}
\setcounter{ex}{0}
\begin{ex}% Câu 1
	Một miếng đồng được khoan bằng máy khoan điện. Máy khoan hoạt động ở công suất $40,0\ W$ và mất $30,0\ s$ để khoan xuyên qua miếng đồng. Nếu $80\%$ năng lượng từ máy khoan làm nóng miếng đồng. Độ tăng nội năng của miếng đồng là bao nhiêu $J$?
	\SA[oly]{$960$}
	\loigiai{
		\begin{itemize}
			\item $ \Delta U = Q + A = \mathcal{P.H}.t + 0 = 40.0,8.30 = 960\ J$.
		\end{itemize}
	}
\end{ex}

\begin{ex}% Câu 2
	Trong hệ thống làm mát của một động cơ, động cơ được làm mát nhờ dòng chất lỏng tuần hoàn đi vào các chi tiết làm mát hấp thu nhiệt và đi ra các ống làm mát để giảm nhiệt độ. Cho rằng nhiệt độ của dòng chất lỏng khi đi ra khỏi các chi tiết cần làm mát là $60^{\circ}C$, chất lỏng này di chuyển qua các ống làm mát (xung quanh ống là $60\ l$ nước ở nhiệt độ $10^{\circ}C$). Sau khi chất lỏng di chuyển qua các ống nhiệt độ giảm xuống còn $30^{\circ}C$. Sau khoảng thời gian $t$ nhiệt độ của nước tăng lên thành $20^{\circ}C$. Biết nhiệt dung riêng của nước và của chất lỏng lần lượt là $1\ cal/g.^{\circ}C$ và $0,5\ cal/g.^{\circ}C$, khối lượng riêng của nước là $\rho = 1000\ kg/m^3$. Khối lượng chất lỏng di chuyển qua ống trong khoảng thời gian $t$ bằng bao nhiêu kg?
	\par \shortans[oly] {$40$}
	
	\loigiai{
		\begin{itemize}
			\item Nhiệt tỏa ra của chất lỏng sẽ được nước hấp thụ: $m_lc_l\Delta t_l = m_nc_n\Delta t_n \Rightarrow m_l.0,5(60-30) = 60.10^3(20-10) \Rightarrow m_l = 40.10^3\ = 40\ kg$.
		\end{itemize}
	}
\end{ex}

\begin{ex}% Câu 3
	Để đóng băng $500\ g$ nước ở $30^\circ C$ thành nước đá, người ta bỏ vào nước trên một khối nước đá ở $-10^\circ C$. Biết nhiệt dung riêng của nước là $c_1 = 4200\ J/kg.K$ và nước đá là $c_2 = 2000\ J/kg.K$; nhiệt nóng chảy riêng của nước đá $\lambda = 3,4.10^5\ J/kg$. Lượng nước đá tối thiểu cần dùng bằng bao nhiêu kg (làm tròn kết quả đến chữ số hàng đơn vị)?
	
	\shortans[oly]{$12$}
	\loigiai{
	\begin{itemize}
		\item $m_n.c_n.\Delta t_n + m_n.\lambda = m_{\text{đ}}.c_{\text{đ}}.\Delta t_{\text{đ}} \Rightarrow 0,5.4200.30 + 0,5.3,4.10^5 = m_{\text{đ}}.2000.10 \Rightarrow m_{\text{đ}} = 11,65\ kg \approx 12\ kg$.
	\end{itemize}
	}
\end{ex}

\begin{ex}% Câu 4
	Biết nhiệt hóa hơi riêng và nhiệt dung riêng của nước lần lượt là $2,256. 10^6\ (J/kg)$ và $4,19.  10^3\ (J/kg.K)$. Khi da người nhận được nhiệt tỏa ra bởi $22,0\ g$ hơi nước lúc đầu ở $100^\circ C$, khi nó được làm lạnh đến nhiệt độ da (khoảng $34^\circ C$). Lượng nhiệt truyền vào da của chúng ta là bao nhiêu $kJ$ (làm tròn kết quả đến chữ số hàng phần mười)? 
	\shortans[oly]{$55,7$}
	\loigiai{
		\begin{itemize}
			\item $Q = m.L + m.c.\Delta t = 55716\ J \approx 55,7\ kJ$.
		\end{itemize}
	}
\end{ex}

\begin{ex}% Câu 5
	Người ta rót vào khối nước đá khối lượng $m_1 = 2\ kg$ một lượng nước $m_2 = 1\ kg$ ở nhiệt độ $t_2 = 10^\circ C$. Khi có cân bằng nhiệt, lượng nước đá tăng thêm $m' = 50\ g$. Sau đó người ta cho hơi nước sôi vào bình trong một thời gian. Sau khi thiết lập cân bằng nhiệt, nhiệt độ của nước là $50^\circ C$. Biết nhiệt dung riêng của nước $c = 4200\ J/kg.K$. Nhiệt nóng chảy riêng của nước đá $\lambda = 3,4.10^5\ J/kg$. Cho nhiệt hoá hơi riêng của nước $L = 2,3.10^6\ J/kg$. Bỏ qua sự trao đổi nhiệt với đồ dùng thí nghiệm. Khối lượng hơi nước đã dẫn vào bằng bao nhiêu $kg$ (làm tròn kết quả đến chữ số hàng phần trăm)?
	\shortans[oly]{$0,53$}
	\loigiai{
		\begin{itemize}
			\item Nhận thấy lượng nước đá tăng thêm $50\ g < 1\ kg$ nên nhiệt độ cân bằng là $0^\circ C$.
			\item Khối lượng nước đá ở $0^\circ C$ là $m_{\text{đ}} = 2,05\ kg$; Khối lượng nước ở $0^\circ C$ là $m_n = 1- 0,05 = 0,95\ kg$.
			\item Phương trình cân bằng nhiệt: $Q_{thu} = Q_{\text{tỏa}} \Rightarrow \lambda m_{\text{đ}} + (m_{\text{đ}}+ m_n).c.\Delta t = L.m_h + m_h.c.\Delta t \Rightarrow 3,4.10^5.2,05+3.4200.50 = 2,3.10^6.m_h+m_h.4200.50 \Rightarrow m_h \approx 0,53\ kg$.
		\end{itemize}
	}
\end{ex}

\begin{ex}% Câu 6
	Để đúc các vật bằng thép, người ta phải nấu chảy thép trong lò. Thép đưa vào lò có nhiệt độ $t_1 = 20^\circ C$, hiệu suất của lò là $60\%$, nghĩa là $60\%$ nhiệt lượng cung cấp cho lò được dùng vào việc đun nóng thép cho đến khi thép nóng chảy. Để cung cấp nhiệt lượng, người ta đã đốt hết $m_t = 200\ kg$ than đá có năng suất toả nhiệt là $q_t = 29.10^6\ J/kg$. Cho biết thép có nhiệt nóng chảy riêng $\lambda = 83,7.10^3\ J/kg$; nhiệt độ nóng chảy là $t_2 = 1400^\circ C$; nhiệt dung riêng ở thể rắn là $c = 0,46\ kJ/kg.K$. Khối lượng của mẻ thép bị nấu chảy bằng bao nhiêu tấn (làm tròn kết quả đến chữ số hàng phần mười)?
	\shortans[oly]{$4,8$}
	\loigiai{
		\begin{itemize}
			\item $Q_{thu} = Q_{\text{tỏa}} \Rightarrow m_t \left( c \Delta t + \lambda \right) = 0,6.m_{than}.q \Rightarrow m(460.1380+83,7.10^3) = 0,6.200.29.10^6 \Rightarrow m \approx 4843\ kg \approx 4,8$ tấn.
		\end{itemize}
	}
\end{ex}
